\documentclass[10pt,a4paper]{amsart}
\usepackage[margin=19mm]{geometry}
\usepackage{amsmath,amssymb,mathtools}
\usepackage[scr=rsfs,cal=euler]{mathalfa}
\usepackage{xfrac}
\usepackage[unicode,hidelinks,bookmarks=false]{hyperref}
\newcommand{\ie}{i.e.\ }
\newcommand{\cf}{cf.\ }
\newcommand{\resp}{respectively\ }
\newcommand{\Subsection}{\subsection}
\providecommand{\nobreakdash}{\nobreak\hbox{-}}
\setlength{\emergencystretch}{2em}
\multlinegap0pt
\usepackage{bm}
\usepackage{bbm}
\usepackage{dsfont}
\usepackage{xspace}
\usepackage{cancel}
\usepackage{mathtools}
\usepackage{amsthm}
\makeatletter
\@ifundefined{equation}{\newtheorem{equation}{Equation}}{}
\@ifundefined{claim}{\newtheorem{claim}[equation]{Claim}}{}
\@ifundefined{lemma}{\newtheorem{lemma}[equation]{Lemma}}{}
\@ifundefined{proposition}{\newtheorem{proposition}[equation]{Proposition}}{}
\makeatother
\makeatletter
\expandafter\ifx\csname taggedthm\endcsname\relax
\newenvironment{taggedthm}[1]
 {\renewcommand\thetaggedthmx{#1}\taggedthmx}
 {\endtaggedthmx}
\fi
\newcommand*{\wt}[1]{\widetilde{#1}}
\makeatother
\title[Dual spaces of geodesic currents]{Dual spaces of geodesic currents}
\author{Luca De Rosa, D\'idac Mart\'inez-Granado}
\date{Source: arXiv:2211.05164v4 (CC BY 4.0)}
\begin{document}
\maketitle

\noindent\emph{Source and licence.} Dual spaces of geodesic currents. Version arXiv:2211.05164v4 (CC BY 4.0), distributed under the Creative Commons Attribution 4.0 International licence, as recorded in the arXiv licence metadata for this exact version and shown on the abstract page https://arxiv.org/abs/2211.05164v4. Full rights notice: licenses/arxiv-2211.05164-CC-BY-4.0.txt.\par\smallskip

\noindent\textit{Editorial scope.} Counted unit: the Proposition whose source label is \texttt{prop:contdistance} in the article (source0.tex lines 1033--1037, source label \texttt{prop:contdistance}), delivered together with the article's own complete original proof of it (source0.tex lines 1038--1093, one \texttt{proof} environment of 56 lines). No mathematics is written, repaired or re-derived: statement and proof are the article's own text line for line. Required context retained uncounted: eq:complements (lines 952--955); lem:fatou (lines 973--984); prop:zeromeasure (lines 1009--1016). Every definition, display and label of those lines is retained unchanged. Omitted from this package and not redistributed: 1--951, 956--972, 985--1008, 1017--1032, 1094--2844. Nothing inside a retained excerpt is omitted or altered: every retained body is delivered exactly as the article's lines are written, so no omission receipt of any kind accompanies this package. Citations: the retained excerpts carry no citation command at all, so no bibliography entry of the article is transcribed and no reference is redistributed. Reference adaptation: none was needed and none was made; every reference inside the delivered unit resolves inside it, so no unresolved reference or citation is produced. Printed slips kept literal: no printed word, symbol, hypothesis or number of the counted statement or of its proof was altered. Layout and class: the article's own class is \texttt{amsart} with packages including geometry, amsmath, latexsym, amssymb, amsthm, graphicx, mathtools, caption, subcaption, xfrac, fix-cm, accents, multicol, color; the wrapper is the lane's pinned \texttt{amsart} editorial wrapper at 10pt on A4, which restates the theorem-like environments the retained text needs and adds the article's own macro definitions that the retained text needs (\texttt{\textbackslash wt}), with extra wrapper packages (bm, bbm, dsfont, xspace, cancel, mathtools). The delivered numbering is the unit's own. Assets: no retained span contains a figure environment, a graphics-inclusion command, a PDF-page inclusion, a table or a TikZ picture; no image file is copied into this package, the package declares no asset dependency and no image file is redistributed with it. Licence: version arXiv:2211.05164v4 of this article is distributed under the Creative Commons Attribution 4.0 International licence, as recorded in the arXiv licence metadata for this exact version and shown on the abstract page https://arxiv.org/abs/2211.05164v4. A full rights notice is delivered at licenses/arxiv-2211.05164-CC-BY-4.0.txt. Characters: every retained excerpt is pure ASCII, so the delivered engine, XeLaTeX with its default Latin Modern text font, reports no missing character.
\begin{equation}
\liminf_n A_n = \left( \limsup_n A_n^c \right)^c.
\label{eq:complements}
\end{equation}

\begin{lemma}
Let $\mu$ be a (non-necessarily finite) measure on a measurable space $X$, and let $(A_n)$ be a sequence of measurable sets.
Then
\[
\liminf_n \mu(A_n) \geq \mu(\liminf_n A_n).
\]
Moreover, if for some $n$, $\mu \left( \cup_{j \geq n} A_j \right)$ is finite, we have
\[
\limsup_n \mu(A_n) \leq \mu(\limsup_n A_n).
\]
\label{lem:fatou}
\end{lemma}

\begin{proposition}
Let $X=\Gamma \backslash \wt{X}$ be a hyperbolic structure. Let $G[\bar{x}]$ denote the set of geodesics through $\bar{x} \in \wt{X}$.
If $\bar{x}$ does not lie on the axis of an element $g \in \Gamma - \{ e\}$, then
\[
\mu(G[\overline{x}])=0.
\]
\label{prop:zeromeasure}
\end{proposition}

\begin{proposition}
If $\mu$ is a geodesic current without atoms, then the function
\[f(\cdot, \cdot) \coloneqq d_{\mu}(\pi_{\mu}(\cdot ), \pi_{\mu}(\cdot)) \colon \wt{X} \times \wt{X} \to \mathbb{R}\] is a continuous function.
\label{prop:contdistance}
\end{proposition}

\begin{proof}
Let $\left( \overline{x_n}, \overline{y_n} \right)$ be a sequence in $\wt{X} \times \wt{X}$ converging to $(\overline{x}, \overline{y})$.
We first show that
\begin{equation} \liminf_n d_{\mu}(\pi_{\mu}(\overline{x_n}), \pi_{\mu}(\overline{y_n})) \geq d_{\mu}(\pi_{\mu}(\overline{x}), \pi_{\mu}(\overline{y}))
\label{eq:lower}
\end{equation}
Note that, since $\mu$ has no atoms, we have, by Proposition~\ref{prop:zeromeasure},
\[
d_{\mu}(\pi_{\mu}(\overline{x_n}), \pi_{\mu}(\overline{y_n})) = \mu(G(\overline{x_n},\overline{y_n})).
\]
for all $n$.
\begin{claim}
If $\gamma \in G(\overline{x},\overline{y})$, then $\gamma \in \cup_{N \geq 1} \cap_{n \geq N} G(\overline{x_n},\overline{y_n})$.
\label{claim:lim_inf}
\end{claim}
\begin{proof}
Since $\gamma \in G(\overline{x},\overline{y})$, $\gamma$ intersects $(\overline{x},\overline{y})$ transversely at some $\overline{z} \in (\overline{x},\overline{y})$. Since
$(\overline{x_n},\overline{y_n})$ is converging in Hausdorff distance to $(\overline{x},\overline{y})$, we have that for $n_0$ large enough, $\gamma$ must intersect transversely all $(\overline{x_n},\overline{y_n})$ for all $n \geq n_0$.
\end{proof}
From Claim~\ref{claim:lim_inf} and Lemma~\ref{lem:fatou}, it follows that 
\[
\mu(G(\overline{x},\overline{y})) \leq \mu \left( \cup_{N \geq 1} \cap_{n \geq N} G(\overline{x_n},\overline{y_n})\right) \leq \liminf_n \mu(G(\overline{x_n},\overline{y_n}))
\]
from which Equation~\ref{eq:lower} follows.
We will now prove 
\begin{equation} \limsup_n d_{\mu}(\pi_{\mu}(\overline{x_n}), \pi_{\mu}(\overline{y_n})) \leq d_{\mu}(\pi_{\mu}(\overline{x}), \pi_{\mu}(\overline{y})).
\label{eq:upper}
\end{equation}




\begin{claim}
If $\gamma$ is a geodesic disjoint from $[\overline{x},\overline{y}]$, then $\gamma \notin \limsup_{n} G(\overline{x_n},\overline{y_n})$.
\label{claim:upper}
\end{claim}
\begin{proof}
Let $A_n=G(\overline{x_n},\overline{y_n})$.
We are assuming $\gamma$ is a geodesic disjoint from $(x,y)$, and we want to show  $\gamma \notin \limsup_{n} A_n$, which, by Equation~\ref{eq:complements} is equivalent to showing
$\gamma \in \liminf_n A_n^c$, i.e., to showing that there exists $N > 0$ so that for all $n \geq N$, $\gamma \notin A_n$.
This follows, since $\gamma$ and $(\overline{x},\overline{y})$ are a definite distance apart, and $(\overline{x_n},\overline{y_n})$ is converging to $(\overline{x},\overline{y})$ in the Hausdorff distance of $\wt{X}$. Thus, for $n_0$ large enough, $(\overline{x_n},\overline{y_n})$ is also disjoint from $\gamma$, for all $n \geq n_0$.
\end{proof}
A geodesic which is not in $G(\overline{x},\overline{y})$ is either disjoint from $(\overline{x},\overline{y})$, or it is the geodesic $\gamma_{x,y}$ determined by the points $x$ and $y$.
Claim~\ref{claim:upper} can be phrased as follows:
\[
\limsup_{n}  A_n \subset G(\overline{x},\overline{y}) - \{ \gamma_{x,y} \}.
\]
Since $\mu$ has no atoms, we have
\[
\limsup_{n}  \mu(A_n ) \leq \mu(\limsup_{n} A_n ) \leq \mu( G(\overline{x},\overline{y}) - \{ \gamma_{x,y} \} ) = \mu( G(\overline{x},\overline{y})),
\]
where we have used the second equation in Lemma~\ref{lem:fatou} in the first inequality, and the last equation plus the monotonicity of $\mu$ in the second inequality.
It remains to explain why the application of Lemma~\ref{lem:fatou} is justified. Since $\overline{x_n} \to \overline{x}$ and $\overline{y_n} \to \overline{y}$, there exists $n_0$ so that for $n \geq n_0$, $(\overline{x_n},\overline{y_n}) \subset B$ for some compact ball $B$. Then \[\cup_{n \geq n_0} A_n \subset G(B).\]
Since $B$ is compact, $\mu(G(B))$ is finite, and thus the application of the Lemma is justified.
This shows Equation~\ref{eq:upper} and finishes the proof.
\end{proof}

\end{document}
